\documentclass[UTF8,fontset=none,11pt,a4paper]{ctexart}
\usepackage{ipara-handbook}

% 新建 Topic / Bridge / Integration 时复制本文件。
% 这里只演示公共视觉组件；认知结构仍执行 handbook_writing_spec.md。
\handbooksetup
  {示例 Handbook：对象、机制与边界}
  {数学一 · 示例 Topic XX}
  {对象 · 机制 · 边界}

\begin{document}

\handbooktitle
  {示例 Handbook：对象、机制与边界}
  {公共样式的最小可编译骨架}
  {Topic XX · LaTeX Template}

\begin{corebox}{母问题}
一个对象为什么需要当前机制？机制改变什么状态，又必须保持什么不变量？
\end{corebox}

\begin{methodbox}{本册定位与 Owner 边界}
正文只定义本册 Own 的机制；相邻机制通过 Use / Bridge / Integration 链接，不为了排版完整而复制第二份定义。
\end{methodbox}

\section{最小生成链}

\[
\boxed{
\text{Problem}
\to \text{Naive Solution}
\to \text{Failure}
\to \text{Mechanism}
\to \text{Invariant}
}
\]

\begin{boundarybox}{概念边界}
\begin{boundarytable}
\boundarytablehead
对象 & ≠ & 表示
  & 对象是被研究的实体；表示是为了计算或推理选择的编码。题目更换坐标而实体不变时，应先追踪不变量
  & 会把表示变化误判成对象变化 \\
\bottomrule
\end{boundarytable}
\end{boundarybox}

\section{流程图}

\begin{handbookdiagram}
\begin{tikzpicture}[node distance=16mm and 20mm]
  \node[ipara node] (input) {输入对象};
  \node[ipara process,right=of input] (mechanism) {机制};
  \node[ipara state,right=of mechanism] (state) {稳定状态};

  \draw[ipara edge] (input) -- node[ipara label,above]{处理} (mechanism);
  \draw[ipara edge] (mechanism) -- node[ipara label,above]{提交} (state);
\end{tikzpicture}
\diagramcaption{箭头标签只承担短事件或动作；原因解释留在正文。}
\end{handbookdiagram}

\section{普通表格}

\begin{handbooktable}{L{2.7cm}YY}
\toprule
\textbf{对象} & \textbf{要问的问题} & \textbf{检查} \\
\midrule
状态 & 当前有哪些合法取值？ & 是否遗漏非法/终止状态？ \\
事件 & 什么触发状态变化？ & 事件方向与箭头语义是否一致？ \\
边界 & 哪些条件不满足时必须停？ & 是否已经交给正确 Owner？ \\
\bottomrule
\end{handbooktable}

\section{代码与标识符}

行内标识符使用 \code{TLB miss}、\code{fork()}；较长路径使用 \codepath{/proc/<pid>/maps}。

\begin{lstlisting}[language=C]
if (condition) {
    update_state();
} else {
    keep_invariant();
}
\end{lstlisting}

\begin{warnbox}{不要用缩小字号掩盖结构问题}
图或表越界时，先缩短标签、调整列职责、重做路由或拆图；整体缩放只是最后手段。
\end{warnbox}

\end{document}
